% !TeX root = ../../Book3.tex
\section{\texorpdfstring{\(\operatorname{Re}s=1\)}{Re(s)=1} 上的无零点性}
\label{b3:sec:2703}

无零点性只需要一个非负三角多项式和局部 Laurent 展开。这里不会证明一条定量宽度的无零区域；所证内容恰是边界直线上每个有限高度没有零点。

\subsection{Hadamard–de la Vallée Poussin 论证}
\label{b3:subsec:270301}

以下经典论证称为\term[hadamard-de-la-vallee-poussin]{Hadamard--de la Vallée Poussin 无零点论证}[Hadamard--de la Vallée Poussin zero-free argument]。设 \(t\ne0\)，若 \(1+\mathrm{i}t\) 是 \(\zeta\) 的 \(m\) 阶零点，则
\(F(s)=-m/(s-1-\mathrm{i}t)+O(1)\) 于该点附近。若 \(1+2\mathrm{i}t\) 也是零点，记阶数为 \(m_2\)，否则记零。

\subsection{正性技巧}
\label{b3:subsec:270302}

恒等式 \(3+4\cos u+\cos2u=2(1+\cos u)^2\ge0\) 给
\begin{equation}\label{b3:eq:27-positivity}
3F(\sigma)+4\operatorname{Re}F(\sigma+\mathrm{i}t)
+\operatorname{Re}F(\sigma+2\mathrm{i}t)\ge0,\quad\sigma>1.
\end{equation}
每个权 \(\Lambda(n)n^{-\sigma}\) 非负，故相乘求和保持不等式。把中间系数改小会丧失排除单零点所需的优势；三个系数的具体数值承担了证明任务。

\subsection{边界估计}
\label{b3:subsec:270303}

令 \(\sigma=1+\varepsilon\)。局部展开给
\(\varepsilon F(1+\varepsilon)\to1\)，
\(\varepsilon\operatorname{Re}F(1+\varepsilon+\mathrm{i}t)\to-m\)，第三项极限为 \(-m_2\)。固定 \(t\) 使每个展开只在一个固定邻域使用，无需关于所有高度的一致估计。

\subsection{零点排除}
\label{b3:subsec:270304}

\begin{theorem}\label{b3:thm:27-zero-free-line}
Riemann zeta 函数在直线 \(\operatorname{Re}s=1\) 上没有零点；\(s=1\) 处仍为单极点。
\end{theorem}
\begin{proof}
若 \(t\ne0\) 处有 \(m\ge1\) 阶零点，将\eqref{b3:eq:27-positivity}乘 \(\varepsilon\) 并取极限，得到 \(3-4m-m_2\ge0\)，但右式不超过负一，矛盾。\(t=0\) 处已知为极点。
\end{proof}
所以 \(F(s)-1/(s-1)\) 在整条闭半平面 \(\operatorname{Re}s\ge1\) 的一个开邻域中全纯；“一个开邻域”允许向左伸出的宽度随高度改变，不暗含统一条带。下一节使用的正是这种局部延拓。

\ProseParagraphBreak

这里可以逐点说明这个邻域如何获得。在不等于一的边界点，\(\zeta\) 全纯且非零，连续性使它在一个足够小的圆盘内仍非零，因而对数导数全纯。在一处，写
\[
 \zeta(s)=\frac{a(s)}{s-1},\qquad a(1)=1,
\]
其中 \(a\) 在小圆盘内全纯非零，则
\[
 F(s)-\frac1{s-1}=-\frac{a'(s)}{a(s)}
\]
在该圆盘全纯。把这些圆盘与原右半平面取并，便得到需要的开邻域。对于每一段有限高度的边界，紧致性还能选出统一的正宽度；当高度增长时，这个宽度却可能缩小。

\ProseParagraphBreak

另一个等价的正性检验从 \(\log\zeta\) 出发。在 \(\sigma>1\) 内逐项应用同一三角多项式并取指数，得
\[
 \zeta(\sigma)^3|\zeta(\sigma+\mathrm it)|^4
              |\zeta(\sigma+2\mathrm it)|\ge1.
\]
若 \(1+\mathrm it\) 是 \(m\ge1\) 阶零点，左侧在 \(\sigma\downarrow1\) 时为
\(O((\sigma-1)^{4m+m_2-3})\to0\)，矛盾。指数 \(4m+m_2-3\) 与前面留数不等式的符号相反，两者实际计算了同一局部阶数。这个版本也再次显示为何需要系数四压过极点贡献三。
